%Packages
\usepackage[utf8]{inputenc}
\usepackage{amsmath,amssymb}
\usepackage{tikz}

%Set the theme and stype
\usetheme{Pittsburgh}
\usecolortheme{rose}
\usefonttheme{serif}
\useinnertheme{rounded}
\useoutertheme{infolines}
\setbeamercovered{dynamic}
%\setbeamertemplate{navigation symbols}{}
\setbeamercovered{invisible}
%For better formating of the tables
\usepackage{booktabs}

%Setting up the parts, lectures and sections
\AtBeginLecture{\subtitle{\insertlecture} \maketitle}
\AtBeginPart{\frame{\partpage}\frame[allowframebreaks]{\tableofcontents[hideallsubsections]}}
\AtBeginSection[]{\frame[allowframebreaks]{\frametitle{Outline - Next Section}\tableofcontents[currentsection,hideothersubsections]}}

%------Change block to subsubsections in the article model-----%
\mode<article>{
    % Redefine the standard Beamer block to become a subsubsection
    \renewenvironment{block}[1]{
        \subsubsection{#1}
    }{
        \par\vspace{0.3cm} % Adds a nice little breathing space after the block
    }
}

%------Change strucutre command to subsubs ection as well in article mode----%
\mode<article>{
    % 1. Redefine \structure to act as an unnumbered subsubsection
    %\renewcommand{\structure}[1]{\subsubsection{#1}}
    \renewcommand{\structure}[1]{\textbf{#1}}

}

%-------Making nice defintion boxes for the reader-------%
\mode<article>{
    \usepackage{tcolorbox}
    
    % 1. "Delete" the boring default article definition
    \let\definition\relax
    \let\enddefinition\relax
    
    % 2. Create a beautiful new Beamer-style box for the reader
    \newtcolorbox{definition}[1][]{
        colback=blue!5!white,      % Very light blue background
        colframe=blue!75!black,    % Dark blue border and header
        fonttitle=\bfseries,       % Bold header text
        title={Definition\ifx\relax#1\relax\else: #1\fi}, % Auto-adds the title
        arc=2mm,                   % Rounded corners
        boxrule=1pt,               % Thickness of the border
        left=2mm, right=2mm, top=2mm, bottom=2mm % Internal padding
    }
}


%-----Setting for including python code----%
\usepackage{listings}
\usepackage{xcolor}

% Define colours for syntax highlighting
\definecolor{codegreen}{rgb}{0,0.6,0}
\definecolor{codegray}{rgb}{0.5,0.5,0.5}
\definecolor{codepurple}{rgb}{0.58,0,0.82}
\definecolor{backcolour}{rgb}{0.95,0.95,0.92}

% Configure the style
\lstdefinestyle{mystyle}{
    backgroundcolor=\color{backcolour},   
    commentstyle=\color{codegreen},
    keywordstyle=\color{magenta},
    numberstyle=\tiny\color{codegray},
    stringstyle=\color{codepurple},
    basicstyle=\ttfamily\footnotesize, % Typewriter font, slightly smaller
    breakatwhitespace=false,         
    breaklines=true,                 
    captionpos=b,                    
    keepspaces=true,                 
    numbers=left,                    
    numbersep=5pt,                  
    showspaces=false,                
    showstringspaces=false,
    showtabs=false,                  
    tabsize=4
}
\lstset{style=mystyle}
%---End python code setting----%

%Define my custom enviroments
\newenvironment<>{recipe}[1][]{%
%  \setbeamercolor{block title}{fg=white,bg=red!75!black}%
  \begin{alertblock}{Recipe (#1)}}{\end{alertblock}}


%Remove the frame title from article and here use just the section and subsection titles.
\mode<article>{\renewcommand{\frametitle}[1]{}}

%\mode<article>{\renewcommand{\lecturetitle}[1]{}}

\definecolor{SolarizedYellow}{RGB}{253, 246, 227} % Body
\definecolor{SolarizedAccent}{RGB}{238, 232, 213} % Title (slightly darker)

\newcommand{\insertcode}[2]{%
  \only<presentation>{%
    \setbeamercolor{block title}{bg=SolarizedAccent, fg=black}
    \setbeamercolor{block body}{bg=SolarizedYellow, fg=black}%
    \begin{block}{#1}
      \lstinputlisting[language=Python]{#2}
    \end{block}%
  }%
  \only<article>{%
    \vspace{0.2cm}
    % Changing 'title' to 'caption' enables automatic numbering
    \lstinputlisting[language=Python, caption={#1}]{#2}
    \vspace{0.2cm}
  }%
}



% Title Formatting
\title{\textbf{MATH19622-1E2: Introduction to Probability and Statistics}\\
\large 4 Lecture Series (90 minutes per Lecture)}
\author{Anthony R. Thornton}
\date{2025/2026}
